% ID: SEL-D0067-OB-E03
% ESTADO: validada. S067.
\textbf{Parte I.} Sean $H$ el préstamo hipotecario y $P$ el personal.

\textbf{a)} La probabilidad de tener alguno es $0{,}60+0{,}50-0{,}20=0{,}90$. Por complemento, la probabilidad de no tener ninguno es $\boxed{0{,}10}$.

\textbf{b)} Tener hipoteca sin préstamo personal tiene probabilidad $0{,}60-0{,}20=0{,}40$. Por tanto, $P(H\mid P^c)=0{,}40/0{,}50=\boxed{0{,}80}$.

\textbf{Parte II.} En el modelo habitual de observaciones independientes, la media muestral es normal de media $110$ y desviación típica $15/\sqrt n$.

\textbf{a)} Con $n=25$, la desviación típica es $3$. El umbral $113$ corresponde a $z=1$. Si $\Phi$ es la distribución normal estándar,
\[
P(\bar X>113)=1-\Phi(1)\approx\boxed{0{,}1587}.
\]
\textbf{b)} Al aumentar $n$, el umbral tipificado $(113-110)/(15/\sqrt n)=\sqrt n/5$ sube. Por tanto, se reduce la probabilidad de cola por encima de él.
